\section{Experiments and Results}
\label{sec:experiments}

\subsection{Corpus-Level Results}

The normalized corpus provides the broadest evidentiary layer in the repository. Table~\ref{tab:dataset-summary} summarizes the headline numbers. Two observations matter immediately. First, the corpus is large enough to support trend and taxonomy analysis rather than anecdotal reasoning. Second, it remains heterogeneous even after normalization. Approximately two-thirds of all incidents are labeled as DeFi, which means a non-trivial fraction of the broader archive still falls outside the paper's focus and must be filtered out before any exploitability claim is made.

\begin{table}[t]
\centering
\caption{Dataset and study summary derived from repository artifacts.}
\label{tab:dataset-summary}
\begin{tabular}{lr}
\toprule
\textbf{Metric} & \textbf{Value} \\
\midrule
Total normalized incidents & 2{,}036 \\
Unique incident identifiers & 2{,}035 \\
DeFi-labeled incidents & 1{,}378 \\
DeFi share of full corpus & 67.68\% \\
Loss parsing coverage within DeFi subset & 68.36\% \\
Deep-audited protocols & 5 \\
Protocols in pure on-chain verification study & 24 \\
Shortlisted replay incidents & 12 \\
Replay cases with concrete logic & 3 \\
\bottomrule
\end{tabular}
\end{table}

Figure~\ref{fig:yearly-incidents} shows the yearly count of DeFi-labeled incidents in the normalized corpus. The curve rises steeply through 2021, peaks in the 2022--2023 period, and then declines in the partial 2025--2026 window. This trend should not be read as a direct measure of protocol insecurity. It also reflects reporting intensity, ecosystem expansion, and dataset timing. Still, it establishes that the corpus is large enough to justify a workflow that prioritizes filtering and mechanism identification over one-off exploit storytelling.

\begin{figure}[t]
\centering
\includegraphics[width=0.93\textwidth]{yearly_defi_incidents.png}
\caption{Yearly DeFi incident counts in the normalized corpus. The trend motivates broad corpus construction but does not by itself imply replayable exploitability.}
\label{fig:yearly-incidents}
\end{figure}

Attack-family counts tell a more useful story than yearly volume. Figure~\ref{fig:family-dist} shows that the normalized DeFi subset is dominated by the \texttt{other} bucket, followed by contract bugs, flash-loan incidents, account compromise, bridge-message abuse, social engineering, and oracle manipulation. The size of the \texttt{other} bucket is itself an important result: public archives remain semantically noisy even after normalization, and any paper that translates those counts directly into exploit conclusions will overstate the prevalence of replayable smart-contract failures.

\begin{figure}[t]
\centering
\includegraphics[width=0.93\textwidth]{incident_family_distribution.png}
\caption{Dominant attack families in the DeFi-labeled corpus. The large residual \texttt{other} category shows why incident counts cannot be treated as a clean exploitability proxy.}
\label{fig:family-dist}
\end{figure}

Loss statistics reinforce the same point. Within the DeFi-labeled subset, the median parsed loss is approximately \$891{,}000, while the top twenty incidents account for roughly \$17.26B in aggregate loss. Yet the appendix shows that many of the largest entries in the dataset fall into broad categories such as \texttt{other} or \texttt{supply\_chain}. This is precisely why the paper separates incident magnitude from replayability. A high-loss event is important evidence for ecosystem risk, but it is not automatically evidence of a replayable pure on-chain contract exploit.

\subsection{Protocol Evidence}

The repository contains two complementary protocol-level datasets. The first consists of five deep audits on Lido, Aave, Uniswap, Curve, and Compound. The second is a 24-protocol pure on-chain verification study that asks a narrower question: which major protocols still expose single-transaction or short-window exploitability under a strict threat model? Figure~\ref{fig:audit-findings} summarizes the manually triaged findings from the audit layer. The 87 findings include 5 critical, 16 high, 42 medium, and 24 low or informational issues. This distribution is useful because it demonstrates that code-level weakness density and replayable exploitability are not identical. Mature protocols still expose many hardening opportunities even when direct pure on-chain exploit paths are scarce.

\begin{figure}[t]
\centering
\includegraphics[width=0.80\textwidth]{audit_findings.png}
\caption{Manually triaged findings across the five deep protocol audits. Scanner output was much larger, but paper claims are based on the triaged set.}
\label{fig:audit-findings}
\end{figure}

The 24-protocol verification study yields an even sharper conclusion. As shown in Figure~\ref{fig:protocol-verdicts}, 20 protocols were judged safe under the study's pure on-chain threat model, 3 were conditional or integration-dependent, and only 1 was confirmed exploitable with high confidence. Table~\ref{tab:protocol-outcomes} makes these categories explicit. The confirmed case is Morpho Blue, where the repository contains an already validated malicious-oracle market proof of concept. The conditional set consists of Pendle, Curve, and BENQI, each of which exposes a real mechanism family but still depends on downstream integrations, empty-market states, or runtime constraints before it becomes a clean exploit.

\begin{figure}[t]
\centering
\includegraphics[width=0.74\textwidth]{protocol_verdicts.png}
\caption{Outcome categories from the 24-protocol pure on-chain verification study. The high-confidence exploitable set is materially smaller than a naive reading of the incident corpus would suggest.}
\label{fig:protocol-verdicts}
\end{figure}

\begin{table}[t]
\centering
\caption{Protocol-study outcomes under the repository's pure on-chain threat model.}
\label{tab:protocol-outcomes}
\begin{tabular}{lp{7.1cm}r}
\toprule
\textbf{Verdict} & \textbf{Interpretation} & \textbf{Count} \\
\midrule
Verified safe & No pure on-chain exploit path identified under the study assumptions & 20 \\
Conditional / integration-dependent & Attack surface exists, but exploitation depends on downstream integration, market state, or additional runtime assumptions & 3 \\
Confirmed exploitable & High-confidence pure on-chain exploit path supported by repository artifacts & 1 \\
\bottomrule
\end{tabular}
\end{table}

The protocol evidence therefore corrects a common intuition. Broad incident feeds make DeFi appear uniformly fragile, whereas protocol-level verification shows concentration: a small number of mechanism families continue to matter, but many high-TVL systems are protected by strong oracle design, constrained architecture, hardened governance, or simply the absence of a direct price-dependent lending path. The result is consistent with the intuition behind robust external price feeds such as Chainlink \citep{chainlinkDocs}, but our contribution is methodological rather than promotional: the repository demonstrates how to convert that intuition into explicit outcome categories.

Expressed as ratios, the result is even clearer. The protocol exposure ratio is approximately $4/24 = 16.7\%$ if conditional and confirmed cases are grouped together, and only $1/24 = 4.2\%$ if the threshold is raised to high-confidence confirmed exploitability. Those two values provide a practical upper and lower bound for how aggressively the current artifact set should be interpreted.

\subsection{Replay Batch Results}

Replay engineering is the final and narrowest evidence layer. The current batch contains twelve shortlisted incidents selected by loss magnitude, exploit-family relevance, and public reproducibility signals. Of those twelve, concrete replay logic has been implemented for three incidents: Euler Finance, BonqDAO and AllianceBlock, and Lendf.Me. Figure~\ref{fig:replay-cases} places those cases alongside UwU Lend, which remains shortlisted but deferred because the exploit path is materially larger and requires a dedicated pass.

\begin{figure}[t]
\centering
\includegraphics[width=0.93\textwidth]{replay_cases.png}
\caption{Replay batch case studies by reported loss. Orange bars indicate incidents with concrete replay logic already encoded in Foundry tests.}
\label{fig:replay-cases}
\end{figure}

\begin{table}[t]
\centering
\caption{Status of the principal replay case studies in the current paper.}
\label{tab:replay-cases}
\begin{adjustbox}{max width=\textwidth}
\begin{tabular}{llllr}
\toprule
\textbf{Case} & \textbf{Family} & \textbf{Reported loss} & \textbf{Status} & \textbf{Fork-verified} \\
\midrule
Euler Finance & Flash-loan-assisted logic abuse & \$197.0M & Logic implemented & No \\
BonqDAO \& AllianceBlock & Oracle manipulation & \$120.0M & Logic implemented & No \\
Lendf.Me & Reentrancy & \$24.7M & Logic implemented & No \\
UwU Lend & Oracle manipulation & \$19.3M & Deferred shortlist & No \\
\bottomrule
\end{tabular}
\end{adjustbox}
\end{table}

These cases are useful precisely because they represent different exploit families. Euler is a flash-loan-assisted liquidation and donation path with logic-level amplification, consistent with post-mortem analyses of the incident \citep{eulerPostmortem}. BonqDAO demonstrates how a manipulated reporting path can collapse collateral assumptions and enable outsized borrowing \citep{bonqIncident}. Lendf.Me remains the canonical callback-driven reentrancy case in which ERC777 token hooks violate assumptions about call order and state finality \citep{lendfmeIncident}. Together, the three cases support a narrower but stronger claim than the incident corpus alone: replayable exploitability in the repository clusters around a small set of mechanism families with well-understood control-flow or oracle preconditions.

The replay conversion rate is currently $3/12 = 25\%$. This number should not be interpreted as a failure. It is evidence that exploit reproduction is the most demanding stage in the workflow. Some incidents lack clean public exploit paths. Some depend on archive RPC state that is unavailable in the current shell configuration. Others, like UwU Lend, are multi-stage enough that forcing them into the batch would reduce correctness. A stronger paper is built by marking those boundaries explicitly rather than pretending that every high-loss incident is equally replayable.
